% Options for packages loaded elsewhere
\PassOptionsToPackage{unicode}{hyperref}
\PassOptionsToPackage{hyphens}{url}
\documentclass[
]{article}
\usepackage{xcolor}
\usepackage{amsmath,amssymb}
\setcounter{secnumdepth}{-\maxdimen} % remove section numbering
\usepackage{iftex}
\ifPDFTeX
  \usepackage[T1]{fontenc}
  \usepackage[utf8]{inputenc}
  \usepackage{textcomp} % provide euro and other symbols
\else % if luatex or xetex
  \usepackage{unicode-math} % this also loads fontspec
  \defaultfontfeatures{Scale=MatchLowercase}
  \defaultfontfeatures[\rmfamily]{Ligatures=TeX,Scale=1}
\fi
\usepackage{lmodern}
\ifPDFTeX\else
  % xetex/luatex font selection
\fi
% Use upquote if available, for straight quotes in verbatim environments
\IfFileExists{upquote.sty}{\usepackage{upquote}}{}
\IfFileExists{microtype.sty}{% use microtype if available
  \usepackage[]{microtype}
  \UseMicrotypeSet[protrusion]{basicmath} % disable protrusion for tt fonts
}{}
\makeatletter
\@ifundefined{KOMAClassName}{% if non-KOMA class
  \IfFileExists{parskip.sty}{%
    \usepackage{parskip}
  }{% else
    \setlength{\parindent}{0pt}
    \setlength{\parskip}{6pt plus 2pt minus 1pt}}
}{% if KOMA class
  \KOMAoptions{parskip=half}}
\makeatother
\usepackage{longtable,booktabs,array}
\newcounter{none} % for unnumbered tables
\usepackage{calc} % for calculating minipage widths
% Correct order of tables after \paragraph or \subparagraph
\usepackage{etoolbox}
\makeatletter
\patchcmd\longtable{\par}{\if@noskipsec\mbox{}\fi\par}{}{}
\makeatother
% Allow footnotes in longtable head/foot
\IfFileExists{footnotehyper.sty}{\usepackage{footnotehyper}}{\usepackage{footnote}}
\makesavenoteenv{longtable}
\setlength{\emergencystretch}{3em} % prevent overfull lines
\providecommand{\tightlist}{%
  \setlength{\itemsep}{0pt}\setlength{\parskip}{0pt}}
\usepackage{bookmark}
\IfFileExists{xurl.sty}{\usepackage{xurl}}{} % add URL line breaks if available
\urlstyle{same}
\hypersetup{
  hidelinks,
  pdfcreator={LaTeX via pandoc}}

\author{}
\date{}

\begin{document}

\section{Invariant rings of ternary cubics and
quartics}\label{invariant-rings-of-ternary-cubics-and-quartics}

\emph{Written and self-checked by GPT-6.1 Sol (OpenAI) in Codex at
Ultra. Original exposition, proofs, exercises and diagrams are CC0.}

Two explicit invariants determine the invariant ring of a ternary cubic.
A ternary quartic needs seven homogeneous parameters and thirteen
algebra generators. We prove these assertions using alternating
contractions, the Hesse pencil, a finite group and exact polynomial
calculations. Determinant differential operators and exterior pairings
provide complementary ways to organize the calculations.

We use the action and small binary covariants from
\href{NOE-INV-01.html}{Covariants, averaging and Hilbert finiteness},
the tensor and symbolic methods from
\href{NOE-INV-02.html}{Polarization, vector invariants and symbolic
evaluation}, and the finite-averaging and Molien formula from
\href{NOE-INV-03.html\#theorem-4-1}{Invariants of finite groups and
Noether's degree bound}. All three precede this lesson. The base field
has characteristic zero; the arguments begin over the complex numbers
and include their descent.

\subsection{1. The complete invariant ring of a ternary
cubic}\label{section-1}

Hilbert's theorem supplies finiteness without choosing generators. For a
ternary cubic, two explicit contractions suffice. The completeness proof
below uses a two-dimensional family and a finite group acting on it. The
counting formula is \href{NOE-INV-03.html\#theorem-4-1}{Molien's
formula, Theorem 4.1 of the preceding lesson}; its finite-averaging
proof is independent of this cubic example.

Write a general cubic with multinomial coefficients: \[
f(x,y,z)=\sum_{i+j+k=3}\frac{3!}{i!j!k!}c_{ijk}x^iy^jz^k.
\] For four independent symbolic rows \(u,v,w,r\), let \(U_f\) replace a
degree-three monomial in each row by the corresponding coefficient
\(c_{ijk}\). Define \protect\phantomsection\label{equation-A1}{}

\[
\begin{gathered}
S(f)=U_f\bigl([uvw][uvr][uwr][vwr]\bigr),\\
H_f=\det\left(\frac{\partial^2f}{\partial x_i\partial x_j}\right),\qquad
T(f)=\left.\frac{d}{ds}S(f+sH_f)\right|_{s=0}.
\end{gathered}
\tag{A1}
\] Every symbolic row in the bracket product occurs three times, so its
evaluation is defined. Equivariance of the umbral map makes \(S\)
special-linear invariant and homogeneous of coefficient degree four. The
chain rule gives \(H_{f\circ h}=\det(h)^2H_f\circ h\). Thus the cubic
\(H_f\), whose coefficients have degree three in those of \(f\), is a
special-linear covariant. Applying invariance of \(S\) to \(f+sH_f\)
shows that \(T\) is invariant, of coefficient degree six. These are
fixed normalizations of the Aronhold invariants; multiplying either by a
nonzero constant changes neither its zero set nor the algebra it
generates.

Consider the \textbf{Hesse pencil of forms}, here including its scalar
parameter: \protect\phantomsection\label{equation-A2}{}

\[
f_{a,b}=aF+bQ,\qquad F=x^3+y^3+z^3,\quad Q=xyz.
\tag{A2}
\] The two parameters \(a,b\) range independently over \(\mathbb C\).

\protect\phantomsection\label{lemma-1-1}{}

\textbf{Lemma 1.1 (restriction to the pencil is injective).} A
polynomial invariant of ternary cubics is determined by its values on
(A2).

\textbf{Proof.} Consider the holomorphic map \[
\operatorname{SL}_3(\mathbb C)\times\mathbb C^2
\longrightarrow \operatorname{Sym}^3(\mathbb C^3)^*,
\qquad (h,a,b)\longmapsto f_{a,b}\circ h.
\] At \((I,1,0)\), its six off-diagonal matrix directions give the six
distinct monomials \(3x_i^2x_j\), \(i\ne j\). Two trace-zero diagonal
directions give \(3(x^3-z^3),3(y^3-z^3)\). The pencil directions give
\(F,Q\). These ten cubics are independent: the mixed monomials and \(Q\)
are distinct, while the two differences and the sum \(F\) span the pure
cubes. The derivative therefore has rank ten. The complex
inverse-function theorem puts a nonempty complex open set in the image.
If an invariant vanishes on the pencil, it vanishes on this image by
invariance, hence on that open set. A polynomial vanishing on a complex
open set is identically zero. \(\square\)

This argument needs no classification of singular cubics, and makes no
claim that every cubic has a Hesse presentation. Density is enough to
determine a polynomial.

We next identify enough symmetries of the pencil to count its possible
invariant polynomials. Fix a primitive cube root \(\zeta\), and put
\protect\phantomsection\label{equation-A3}{}

\[
d=3(\zeta^2-\zeta),\qquad d^2=-27,
\qquad
A=\begin{pmatrix}\zeta&0\\0&1\end{pmatrix},\qquad
B=\frac1d\begin{pmatrix}3&1\\18&-3\end{pmatrix}.
\tag{A3}
\] These matrices act on the column \((a,b)^t\). Both actions arise from
determinant-one substitutions on \((x,y,z)\). For \(A\), choose
\(\rho^3=\zeta\) and substitute
\(\rho\operatorname{diag}(\zeta^{-1},1,1)(x,y,z)^t\). Its determinant is
one, and it sends \(aF+bQ\) to \(\zeta aF+bQ\).

For \(B\), let \[
M=\begin{pmatrix}
1&1&1\\1&\zeta&\zeta^2\\1&\zeta^2&\zeta
\end{pmatrix}.
\] Its determinant is \(d\). Direct multiplication gives \[
F(Mx)=3F(x)+18Q(x),\qquad Q(Mx)=F(x)-3Q(x).
\] Choose \(\sigma^3=d^{-1}\). Substitution by \(\sigma M\) has
determinant one and induces exactly \(B\). Consequently, restriction of
an invariant of ternary cubics is invariant under the group
\(\Gamma=\langle A,B\rangle\).

\protect\phantomsection\label{lemma-1-2}{}

\textbf{Lemma 1.2 (the finite pencil group).} The group \(\Gamma\) has
24 elements, and \protect\phantomsection\label{equation-A4}{}

\[
\sum_{m\geq0}\dim\bigl(\mathbb C [a,b]^\Gamma_m\bigr)t^m
=\frac1{(1-t^4)(1-t^6)}.
\tag{A4}
\]

\textbf{Proof.} Put \(C=ABA^{-1}\) and \(D=A^2BA^{-2}\). Multiplying the
displayed two-by-two matrices, using \(1+\zeta+\zeta^2=0\), gives \[
A^3=I,\qquad B^2=C^2=D^2=-I,\qquad
BC=D=-CB,
\] and conjugation by \(A\) cycles \(B,C,D\). In particular
\(\{I,-I,\pm B,\pm C,\pm D\}\) is a quaternion subgroup of eight
distinct matrices. The three cosets obtained by multiplying it by
\(I,A,A^2\) form a group containing both generators. They are distinct
because their determinants are respectively \(1,\zeta,\zeta^2\). This
proves the order 24 and gives every element without an appeal to a group
classification.

The same multiplication gives this entire trace and determinant list:

{\def\LTcaptype{none} % do not increment counter
\begin{longtable}[]{@{}rll@{}}
\toprule\noalign{}
Number of matrices & Trace & Determinant \\
\midrule\noalign{}
\endhead
\bottomrule\noalign{}
\endlastfoot
1 & \(2\) & \(1\) \\
1 & \(-2\) & \(1\) \\
6 & \(0\) & \(1\) \\
4 & \(1+\zeta\) & \(\zeta\) \\
4 & \(-1-\zeta\) & \(\zeta\) \\
4 & \(\zeta\) & \(\zeta^2\) \\
4 & \(-\zeta\) & \(\zeta^2\) \\
\end{longtable}
}

For a two-by-two matrix the Molien denominator is
\(1-\operatorname{tr}(g)t+\det(g)t^2\). Averaging over \(g\) or
\(g^{-1}\) gives the same sum, so the convention for the induced
function action does not change it. Pairing the last four rows and using
\(\zeta+\zeta^2=-1\) gives \[
\frac1{24}\left(
\frac1{(1-t)^2}+\frac1{(1+t)^2}+\frac6{1+t^2}
+\frac{4(2+t)}{(1-t)(1+t+t^2)}
+\frac{4(2-t)}{(1+t)(1-t+t^2)}\right).
\] Multiplication by \((1-t^4)(1-t^6)\) reduces the parenthesis to 24;
this is (A4). Molien's formula is an identity of formal power series, so
this calculation determines every degree. \(\square\)

The restrictions of our two global invariants can now be calculated.
Their normalizations are \protect\phantomsection\label{equation-A5}{}

\[
\begin{split}
S(f_{a,b})&=\frac{b^4-216a^3b}{54},\\
T(f_{a,b})&=\frac4{27}
\bigl(b^6+540a^3b^3-5832a^6\bigr).
\end{split}
\tag{A5}
\] Here is a small way to check all constants. Invariance under \(A\)
and degree four leave only \(a^3b,b^4\) in the first restriction. In the
product of four determinants in (A1), the coefficient of the monomial
using each of the three coordinates of each row once is 24. The
coefficient using a pure cube in each of three rows and the product of
coordinates in the remaining row, summed over the row choices, is
\(-24\). Since \(c_{111}=b/6\) and \(c_{300}=c_{030}=c_{003}=a\) on the
pencil, these give \(24(b/6)^4-24a^3(b/6)\), the first formula. Both
coefficients follow by expanding the determinants as signed sums of
their six permutations; the symbolic construction specifies the
calculation completely.

For the second formula, the Hessian matrix has diagonal entries
\(6ax,6ay,6az\) and off-diagonal entries \(bz,by,bx\). Its determinant
is \[
H_{f_{a,b}}=-6ab^2F+(216a^3+2b^3)Q.
\] Differentiate the first formula in the parameter direction
\((-6ab^2,216a^3+2b^3)\). This gives exactly the second formula of (A5).
Finally their parameter Jacobian is
\protect\phantomsection\label{equation-A6}{}

\[
\det\frac{\partial(S(f_{a,b}),T(f_{a,b}))}{\partial(a,b)}
=-\frac{256}{9}a^2(b^3+27a^3)^2,
\tag{A6}
\] which is nonzero as a polynomial. Thus these two restrictions are
algebraically independent. For completeness, in characteristic zero a
polynomial relation of least positive degree would, on differentiation
and inversion of this nonzero Jacobian over \(\mathbb C(a,b)\), give two
relations of smaller degree. Its two partial derivatives must therefore
be zero polynomials, making it constant, a contradiction.

\protect\phantomsection\label{theorem-1-3}{}

\textbf{Theorem 1.3 (Aronhold's generators).} The invariant ring of a
ternary cubic is \protect\phantomsection\label{equation-A7}{}

\[
\mathbb C [\operatorname{Sym}^3(\mathbb C^3)^*]^{\operatorname{SL}_3}
=\mathbb C [S,T],\qquad \deg S=4,\quad\deg T=6.
\tag{A7}
\] There is no algebraic relation between the two generators.

\textbf{Proof.} Algebraic independence makes the subalgebra of
\(\mathbb C [a,b]^\Gamma\) generated by the two restrictions a
polynomial algebra, with series \(1/((1-t^4)(1-t^6))\). Lemma 1.2 gives
exactly that series for the whole invariant ring of the pencil. The
inclusion is graded, and each degree is finite-dimensional. Equality of
its dimensions in every degree therefore proves equality of the two
algebras.

Now let \(P\) be any invariant of ternary cubics. Its restriction is a
\(\Gamma\)-invariant, hence a polynomial \(q(S(f_{a,b}),T(f_{a,b}))\).
Both \(P\) and \(q(S,T)\) are global invariants with the same
restriction. Lemma 1.1 makes them equal. Independence of the
restrictions also rules out a global relation. \(\square\)

In particular, \(S(F)=0,T(F)=-864\), whereas
\(S(xyz)=1/54,T(xyz)=4/27\). A cubic with zero Hessian need not be
recognized by a smooth normal form; the polynomial theorem applies to
all coefficient values, including singular and zero cubics. Its nullcone
is exactly \(V(S,T)\).

The construction and completeness statement are due to Aronhold. The
Hesse pencil and its finite symmetry group belong to the classical
geometry surveyed by Artebani and Dolgachev in
\href{https://arxiv.org/abs/math/0611590}{The Hesse pencil of plane
cubic curves}. The proof here gives the explicit substitutions, finite
group, generating invariants and all-degree dimension comparison, rather
than requiring their proofs as an external input.

The same theorem holds over any field of characteristic zero. Indeed the
constructions have rational coefficients. In a fixed coefficient degree,
invariance is the common kernel of the six rational matrices induced by
\(x_j\partial/\partial x_i\), \(i\ne j\). Annihilation implies
invariance under their unipotent subgroups by their finite exponential
series, and these subgroups generate the special linear group by
elementary matrix operations. Taking a kernel of these finite matrices
commutes with extension of the ground field. The monomials \(S^iT^j\)
form a basis of the kernel over \(\mathbb C\) by (A7); rational linear
independence and rank then make them a basis over every
characteristic-zero field. This descent does not assume that the entire
field embeds into \(\mathbb C\).

For the determinant calculations in Sections 2--3, let \(W\) be a
polynomial representation of \(G=\operatorname{SL}_n(\mathbb C)\), put
\(A=\mathbb C[W]\), and let \(G\) act by \((gF)(w)=F(g^{-1}w)\). These
are the conventions of Sections 8--11 of the first lesson.

\subsection{2. Cayley's determinant differential
operator}\label{section-2}

Let \(X=(x_{ij})\), \(\Delta_X=\det X\), and \[
\Omega_X=\det(\partial_{x_{ij}})
=\sum_{\sigma\in S_n}\operatorname{sgn}(\sigma)
\prod_{i=1}^n\partial_{x_{i,\sigma(i)}}.
\]

\protect\phantomsection\label{theorem-2-1}{}

\textbf{Theorem 2.1 (Cayley's identity).} For every positive integer
\(s\), \protect\phantomsection\label{equation-H5}{}

\[
\Omega_X\Delta_X^s
=s(s+1)\cdots(s+n-1)\Delta_X^{s-1}.
\tag{H5}
\]

\textbf{Proof.} The chain rule for left and right matrix multiplication
shows that \(\Omega_X\Delta_X^s\) transforms, on either side, with
determinant exponent \(s-1\). On invertible matrices the action is
transitive, so the result is a scalar times \(\Delta_X^{s-1}\);
polynomial density makes this true on all matrices. Evaluate the scalar
at \(X=I\).

For a fixed permutation \(\sigma\), the derivative
\(\prod_i\partial_{x_{i,\sigma(i)}}\Delta_X^s\) at \(I\) is the
coefficient of \(\prod_i t_{i,\sigma(i)}\) in \(\det(I+T)^s\). In formal
power series, \[
\det(I+T)^s
=\exp\left(s\sum_{r\geq1}\frac{(-1)^{r-1}}r
\operatorname{tr}(T^r)\right).
\] To obtain that square-free monomial, each term in a trace must trace
one cycle of \(\sigma\): every selected entry has its prescribed
outgoing and incoming index, so a closed path cannot leave its cycle. A
cycle of length \(r\) occurs \(r\) times in \(\operatorname{tr}(T^r)\),
canceling the denominator \(r\). It contributes \(s(-1)^{r-1}\). The
exponential joins disjoint cycles, with its factorial canceling their
orderings. The coefficient is therefore
\(\operatorname{sgn}(\sigma)s^{c(\sigma)}\), where \(c(\sigma)\) counts
cycles, including fixed points.

Summing with the signs in \(\Omega_X\) gives
\(\sum_\sigma s^{c(\sigma)}\). This sum equals \(s(s+1)\cdots(s+n-1)\):
when inserting the last letter into a permutation of \(n-1\) letters, it
either forms a singleton, giving a factor \(s\), or is inserted after
one of the \(n-1\) existing letters in a cycle, preserving the number of
cycles. The recurrence starts with the value \(s\) for \(n=1\). This
proves (H5). \(\square\)

There is an algebraic form of compact averaging. For a homogeneous
matrix polynomial \(P(X)\) of degree \(rn\), put
\protect\phantomsection\label{equation-H6}{}

\[
L(P)=\frac{(\Omega_X^rP)(0)}{c_r},\qquad
c_r=\prod_{s=1}^r s(s+1)\cdots(s+n-1);
\quad c_0=1.
\tag{H6}
\] Put \(L(P)=0\) on homogeneous degrees not divisible by \(n\), and
extend linearly.

\protect\phantomsection\label{proposition-2-2}{}

\textbf{Proposition 2.2 (the Omega projection).} \(L\) descends to a
bi-invariant linear functional on
\(\mathbb C [G]=\mathbb C [x_{ij}]/(\Delta_X-1)\), with \(L(1)=1\).
Applying \(L\) to the matrix coefficients of the action on \(A\) gives
the Reynolds operator of \href{NOE-INV-01.html\#theorem-9-1}{Theorem
9.1}.

\textbf{Proof.} The determinant chain rule makes \(L\) invariant under
left and right multiplication by \(G\). By
\href{NOE-INV-01.html\#equation-H1}{(H1)}, the left-invariant
polynomials of degree \(rn\) are exactly the scalar multiples of
\(\Delta_X^r\); degrees not divisible by \(n\) have none. Complete
reducibility from \href{NOE-INV-01.html\#theorem-8-3}{Theorem 8.3}
therefore implies that an invariant functional on each homogeneous piece
is determined by its value on this single invariant vector, or is zero
when no such vector exists. Cayley's identity gives \(L(\Delta_X^r)=1\).

For homogeneous \(P\) of degree \(rn\), the two functionals
\(P\mapsto L(P)\) and \(P\mapsto L(\Delta_XP)\) are invariant and both
take the value one on \(\Delta_X^r\). Hence they coincide by that
uniqueness. For other degrees both are zero. Thus \(L((\Delta_X-1)P)=0\)
for every polynomial \(P\), proving descent.

For \(F\in A\), its translate \(gF\) is a polynomial in the entries of
\(g\) with coefficients in a finite-dimensional part of \(A\). Choose
any representative \(P(X)\) for those matrix coefficients in the
quotient coordinate ring and apply \(L\) coefficientwise. Descent makes
the answer independent of that choice. Bi-invariance makes the resulting
map equivariant with invariant image. Since \(L(1)=1\), it is the
identity on invariants and hence a projection. It preserves the degree
in \(W\). Uniqueness in \href{NOE-INV-01.html\#theorem-9-1}{Theorem 9.1}
identifies it with \(\rho\). \(\square\)

For \(n=1\), these formulas are evaluation at \(X=1\), as they should be
for the trivial group.

\subsection{3. Decomposing a matrix product}\label{section-3}

Noether's papers \emph{On the invariant theory of forms in several
variables} replace a collection of determinant expressions by a matrix
product that retains entire exterior-power rows. Here is one of the
decomposition identities in modern notation, with its signs and its
reduction step.

Let \(\alpha_1,\ldots,\alpha_\rho,\beta_1,\ldots,\beta_\sigma\) be
covectors on \(\mathbb C^n\), and let \(y_1,\ldots,y_{\rho+\sigma}\) be
vectors, where \(\rho+\sigma\leq n\). Define the matrix product by \[
\langle\alpha_1\wedge\cdots\wedge\alpha_\rho,\,
y_1\wedge\cdots\wedge y_\rho\rangle
=\det(\alpha_i(y_j))_{1\leq i,j\leq\rho}.
\] It depends on the exterior-power rows, not on the individual factors
chosen to represent them.

\protect\phantomsection\label{proposition-3-1}{}

\textbf{Proposition 3.1 (decomposition identity).} Put
\(q=\alpha_1\wedge\cdots\wedge\alpha_\rho\) and
\(a=\beta_1\wedge\cdots\wedge\beta_\sigma\). Then
\protect\phantomsection\label{equation-H7}{}

\[
\langle q\wedge a,y_1\wedge\cdots\wedge y_{\rho+\sigma}\rangle
=\sum_{\substack{I\subset\{1,\ldots,\rho+\sigma\}\\|I|=\rho}}
(-1)^{\sum_{i\in I}i-\rho(\rho+1)/2}
\langle q,y_I\rangle\langle a,y_{I^c}\rangle.
\tag{H7}
\] The factors within \(y_I,y_{I^c}\) are in increasing index order.

\textbf{Proof.} Form the \((\rho+\sigma)\)-square pairing matrix with
its \(\alpha\)-rows first and its \(\beta\)-rows last. Expand its
determinant along the first \(\rho\) rows. Choosing their columns
\(I=\{i_1<\cdots<i_\rho\}\) leaves columns \(I^c\) for the other rows.
The column shuffle has sign \((-1)^{\sum_j(i_j-j)}\), the displayed
exponent. The two resulting minors are the two pairings on the right.
This proves (H7). Since both factors are exterior pairings, the formula
also proves that the apparent dependence on the individual vectors
enters only through the exterior product on the left. \(\square\)

This is Noether's 1911 paper, Section 5, equation (32). Her exponent
adds rather than subtracts \(\rho(\rho+1)/2\); the difference is the
even integer \(\rho(\rho+1)\), so the signs agree. Her 1910 announcement
describes why retaining the exterior-power rows makes the symbolic
representation explicit.

The case \(\rho=1,\sigma=r-1\) is also a useful reduction formula:
\protect\phantomsection\label{equation-H8}{}

\[
\langle\alpha\wedge b,y_1\wedge\cdots\wedge y_r\rangle
=\sum_{j=1}^r(-1)^{j-1}
\alpha(y_j)\langle b,y_1\wedge\cdots\widehat{y_j}\cdots\wedge y_r\rangle.
\tag{H8}
\] It is the first-row Laplace expansion, corresponding to the product
identity of Section 3, equation (16) of the same paper. It replaces a
pairing involving \(r\) rows by products involving one and \(r-1\) rows.
Iterating reduces it to products of ordinary scalar pairings, with the
alternating signs preserving the exterior row.

For example, in dimension three it gives the entire reduction
\protect\phantomsection\label{equation-H9}{}

\[
\begin{split}
[\alpha\beta\gamma]\,[y_1y_2y_3]
={}&\alpha(y_1)\bigl(\beta(y_2)\gamma(y_3)-\beta(y_3)\gamma(y_2)\bigr)\\
&-\alpha(y_2)\bigl(\beta(y_1)\gamma(y_3)-\beta(y_3)\gamma(y_1)\bigr)\\
&+\alpha(y_3)\bigl(\beta(y_1)\gamma(y_2)-\beta(y_2)\gamma(y_1)\bigr).
\end{split}
\tag{H9}
\] The left side is the determinant of the pairing matrix by
multiplicativity of determinants; (H8) gives exactly the right side.
Every term is an invariant. If its covector rows are symbolic letters
assigned to forms, umbral evaluation preserves the equality by linearity
and gives a relation among actual concomitants. This is a concrete
reduction inside invariant expressions, with no non-invariant
coefficient introduced.

\subsection{4. Ternary forms and the scope of Noether's
computation}\label{section-4}

For a ternary form, symbolic letters have three coordinates, and
brackets are \([abc]\). A bracket consumes one degree from each of three
letters. For a quartic letter every symbolic expression must still use
exactly four degrees of that letter before evaluation. The scalar-weight
calculation in Section 1 of the first lesson gives \[
3w=4\delta-\omega
\] for a ternary quartic covariant of coefficient degree \(\delta\) and
vector-variable order \(\omega\). A concomitant with dual-variable order
\(\nu\) instead has \(3w=4\delta-\omega+\nu\), since the dual variables
scale inversely. These formulas refer to our inverse-substitution
convention.

For example, the ternary Hessian \(\det(f_{x_ix_j})\) has coefficient
degree three and order six, so weight two. The chain rule gives its
covariance directly: the Hessian matrix after inverse substitution is
\(g^{-\mathsf t}\operatorname{Hess}(f)(g^{-1}x)g^{-1}\); taking its
determinant gives the factor \(\det(g)^{-2}\). The cubic-degree quartic
invariant \[
U([abc]^4)
\] uses degree four in each of three letters and has weight four. Its
nonvanishing can be checked on \(x^4+y^4+z^4\): each letter must use a
single coordinate to degree four, so the six determinant-permutation
terms survive, each with positive sign, giving value six.

\subsubsection{A ternary quadratic}\label{a-ternary-quadratic}

For a quadratic \(f=x^{\mathsf t}Qx\) with symmetric \(3\times3\) matrix
\(Q\), \(\det Q\) is an invariant of coefficient degree three and weight
two. It generates the invariant ring.

\textbf{Proof.} Over an algebraic closure, diagonalize a nondegenerate
symmetric form by congruence: completing a square splits off a nonzero
diagonal entry, and induction diagonalizes its orthogonal complement. A
nonzero diagonal entry can be arranged whenever the form is nonzero,
because the associated bilinear form is recovered from quadratic values
by polarization and \(2\ne0\). Taking square roots of the nonzero
diagonal entries shows that two nondegenerate forms are general-linearly
congruent. If their determinants agree, a congruence between them has
determinant \(1\) or \(-1\). In the latter case compose it with a
reflection isometry of one diagonalized form, which has determinant
\(-1\); the congruence is then special-linear.

For an invariant polynomial \(F(Q)\), put
\(h(t)=F(\operatorname{diag}(t,1,1))\). This is a polynomial. The orbit
assertion gives \(F(Q)=h(\det Q)\) on the dense nondegenerate locus, so
the equality is global. Descent to \(k\) follows by evaluating \(h\) on
that same diagonal slice; its coefficients belong to \(k\). Thus the
invariant ring is \(k [\det Q]\). \(\square\)

\subsubsection{Relative completeness in the ternary
quartic}\label{relative-completeness-in-the-ternary-quartic}

Noether's 1907 note and 1908 dissertation study a system containing
covariants, contravariants and mixed forms of a general ternary quartic.
The original term ``biquadratic'' here means degree four. Her count is
\textbf{331 forms and subsystems, complete relative to two quadratic
forms}. It is not the size of an absolutely complete covariant system,
and it is not a list of 331 generators of the scalar invariant ring.

The distinction is explicit in the 1908 introduction, printed page 25:
the relatively complete system is formed modulo the pair \((\rho,t)\),
and transvection with the system of that pair to obtain an absolutely
complete system is left for further work. The closing summary on printed
page 89 again states relative completeness for 331 forms and subsystems.
The 1907 note describes the same termination at a pair of quadratic
forms.

A relative system supplies reduction expressions with additional moduli
still present. Completing the calculation requires the simultaneous
system of those moduli and the relevant foldings with the relative
system. Merely counting the surviving forms before that last step cannot
establish absolute completeness. The distinction is already visible in
our elementary coefficient proof: computations after inverting \(a\)
gave the localized kernel; the separate saturation argument was needed
to obtain the full polynomial kernel.

\subsubsection{Modern scalar invariants}\label{modern-scalar-invariants}

The scalar invariant ring of ternary quartics has seven homogeneous
parameters of degrees \[
3,6,9,12,15,18,27,
\] and thirteen algebra generators of degrees \[
3,6,9,9,12,12,15,15,18,18,21,21,27.
\] These are different assertions. A homogeneous system of parameters
gives a polynomial subalgebra over which the ring is finite; it need not
itself generate the algebra.

Dixmier proved the parameter assertion in \emph{On the projective
invariants of quartic plane curves} (1987). Ohno supplied the
thirteen-generator result, and Elsenhans gave an explicit computational
verification and another generator system in \emph{Explicit computations
of invariants of plane quartic curves}, Sections 5--6 (2015). His
verification computes lower ranks by evaluations modulo a prime and
compares them with the Hilbert-series dimensions through a proved degree
bound.
\href{https://math.uni-paderborn.de/fileadmin-eim/mathematik/AG-Computeralgebra/Publications-elsenhans/inv_pl_quart_5a.pdf}{Elsenhans's
primary paper}

Sections 5--6 prove the seven-parameter and thirteen-generator theorems,
including the exact all-degree and finite independence calculations.

\subsection{5. Seven parameters for a ternary quartic}\label{section-5}

We now prove the seven-parameter assertion, including the finite
calculations on which it depends. This is Dixmier's theorem. The
singular-form reduction also appears in Lercier, Liu, Lorenzo García and
Ritzenthaler, \href{https://arxiv.org/pdf/1803.05816v5}{\emph{Reduction
type of smooth plane quartics}, Theorem 4.1}. The proof uses explicit
contractions, polynomial identities and a rational Jacobian certificate.

Write a quartic with \textbf{ordinary}, rather than multinomially
normalized, coefficients: \[
F(x,y,z)=\sum_{i+j+k=4}c_{ijk}x^iy^jz^k.
\] Put \(R=k [c_{ijk}]^{\operatorname{SL}_3}\). A \textbf{homogeneous
system of parameters} here means algebraically independent homogeneous
invariants \(q_1,\ldots,q_7\) such that \(R\) is a finite module over
\(k [q_1,\ldots,q_7]\). Establishing independence alone would not
establish this conclusion.

\subsubsection{The contractions defining six
parameters}\label{the-contractions-defining-six-parameters}

For forms occupying dual variable spaces, write
\(A\mathbin{\lrcorner}B=A(\partial)B\). For example, \[
\left(\sum a_{ijk}x^iy^jz^k\right)\mathbin{\lrcorner}B(u,v,w)
=\sum a_{ijk}\partial_u^i\partial_v^j\partial_w^kB.
\] This is an equivariant contraction: the chain rule cancels the
transformations of the variables against those of the derivatives.

Take vectors \(r,s\), and restrict \(F\) to their plane: \[
F(rX+sY)=aX^4+bX^3Y+cX^2Y^2+dXY^3+eY^4.
\] The binary quartic invariants in our earlier normalization are \[
i=ae-\frac{bd}{4}+\frac{c^2}{12},\qquad
j=\frac{ace}{6}+\frac{bcd}{48}-\frac{ad^2+b^2e}{16}-\frac{c^3}{216}.
\] Define the ternary contravariants \(\sigma\) and \(\psi\) by
\protect\phantomsection\label{equation-B25}{}

\[
\sigma(r\times s)=i(F(rX+sY)),\qquad
\psi(r\times s)=j(F(rX+sY)).
\tag{B25}
\] They have coefficient degrees two and three, and dual-variable orders
four and six.

Here is why these definitions give polynomials. Two plane bases with the
same nonzero cross product differ by a determinant-one binary change of
basis, so their values of \(i,j\) agree. Scaling the cross product
scales them to the fourth and sixth powers. On the chart \(u\ne0\), use
\(r=(-v,u,0)\), \(s=(-w,0,u)\); their cross product is \(u(u,v,w)\).
Thus \protect\phantomsection\label{equation-B26}{}

\[
\begin{gathered}
\sigma(u,v,w)=\frac{i(F(-vX-wY,uX,uY))}{u^4},\\
\psi(u,v,w)=\frac{j(F(-vX-wY,uX,uY))}{u^6}.
\end{gathered}
\tag{B26}
\] The corresponding constructions on \(v\ne0\) and \(w\ne0\) agree on
overlaps. A denominator in the first chart could only have the
irreducible factor \(u\). Regularity on the second chart excludes a pole
at the generic point of \(u=0\), so this factor cancels. This proves
that (B26) is polynomial in the coefficients and dual variables.
Homogeneity gives the stated orders. Applying
\(g\in\operatorname{SL}_3\) to a plane basis sends its cross product to
\(g^{-\mathsf t}(r\times s)\); hence (B25) proves contravariance.

Set \protect\phantomsection\label{equation-B27}{}

\[
\rho=\frac{F\mathbin{\lrcorner}\psi}{144},\qquad
\tau=\frac{\rho\mathbin{\lrcorner}F}{12}.
\tag{B27}
\] The first is a quadratic contravariant of coefficient degree four;
the second is a quadratic covariant of degree five. Write
\(\rho=u^{\mathsf t}Au\), \(\tau=x^{\mathsf t}Bx\), using symmetric
matrices with half the mixed-term coefficient in each off-diagonal
entry. Define \protect\phantomsection\label{equation-B28}{}

\[
\begin{aligned}
I_3&=\frac{F\mathbin{\lrcorner}\sigma}{72},&
I_9&=\operatorname{tr}(BA),\\
I_{12}&=\det A,& I_{15}&=\det B,\\
I_{18}&=\operatorname{tr}(\operatorname{adj}B\operatorname{adj}A).
\end{aligned}
\tag{B28}
\] For \(I_6\), put \(a_\alpha=c_\alpha\alpha!/4!\), and index a
six-by-six matrix by the degree-two monomials:
\protect\phantomsection\label{equation-B29}{}

\[
I_6=\det(a_{\alpha+\beta})_{|\alpha|=|\beta|=2}.
\tag{B29}
\] This is the catalecticant determinant. In particular, its entries use
normalized coefficients even though \(F\) was written with ordinary
coefficients.

All six expressions are invariants. Under \(g\in\operatorname{SL}_3\),
the quadratic matrices become \(A\mapsto gAg^{\mathsf t}\) and
\(B\mapsto g^{-\mathsf t}Bg^{-1}\), so the traces, determinants and
adjugate trace in (B28) are unchanged. The catalecticant is a bilinear
form on a symmetric-square space; its determinant changes by the square
of \(\det(\operatorname{Sym}^2g)=\det(g)^4=1\). The identity for this
last determinant follows first on diagonal matrices, where each diagonal
entry occurs with total exponent four, and then on elementary matrices,
whose symmetric-square matrices are unipotent. These generate the
special linear group. The remaining contraction in (B28) is equivariant
and has order zero. Their coefficient degrees are precisely their
subscripts. On \(F=x^4+y^4+z^4\), \(I_3=1\), fixing the normalization.

\subsubsection{A degree-27 discriminant, with a determinant
formula}\label{a-degree-27-discriminant-with-a-determinant-formula}

We need both the zero set and a computable formula for the seventh
parameter. We give the requisite resultant argument explicitly.

\protect\phantomsection\label{lemma-5-1}{}

\textbf{Lemma 5.1 (three cubics).} There is a polynomial
\(\operatorname{Res}(g_1,g_2,g_3)\), of degree nine in each cubic's
coefficients, whose vanishing is equivalent to the three homogeneous
cubics having a common point of \(\mathbb P^2\). Normalize it to take
value one on \((x^3,y^3,z^3)\). It is a determinant quotient of the
degree-seven part of the cubic Koszul complex.

\textbf{Proof.} Work first over \(\mathbb C\), and let \(S_d\) denote
the space of ternary forms of degree \(d\). We spell out the elementary
dimension facts used here. Dimension of an irreducible coefficient
variety means the transcendence degree of its function field. An image
parametrized by \(r\) independent coordinates has dimension at most
\(r\), because its coordinate functions lie in that rational function
field. A Jacobian of rank \(r\) gives the reverse bound: differentiating
a polynomial relation of least degree would otherwise give a nonzero
relation of smaller degree.

An irreducible closed subset of \(\mathbb A^n\) of dimension \(n-1\) has
a single irreducible defining polynomial. To see this, take an
irreducible factor \(q\) of a nonzero polynomial in its prime ideal. A
linear coordinate change makes \(q\) monic in the last coordinate:
choose that coordinate's direction where the highest homogeneous part of
\(q\) is nonzero. If the prime ideal contained \(h\notin(q)\), the
nonzero Sylvester resultant of \(q,h\) in the last coordinate would give
a relation among the other \(n-1\) coordinates. It is nonzero by Gauss's
lemma, since \(q\) is irreducible and monic; the Sylvester adjugate
gives an identity expressing it as a combination of \(q,h\). The other
coordinates would then have transcendence degree at most \(n-2\), and
the last is algebraic over them by the monic equation \(q\). This
contradicts dimension \(n-1\). Thus the prime ideal is \((q)\). These
observations justify the parameter counts and hypersurface assertions
below.

The incidence set \[
Z=\{(g_1,g_2,g_3,P):g_1(P)=g_2(P)=g_3(P)=0\}
\] is irreducible of dimension 29. Indeed, over each point of
\(\mathbb P^2\) the three evaluations impose three independent linear
equations on the 30 coefficients. On a coordinate chart this is a
trivial affine bundle of dimension 27 over a two-dimensional chart, and
the charts have dense overlaps.

Its image in coefficient space is closed. To check this without assuming
an elimination theorem, a triple outside the image has affine common
zero set consisting only of the origin. The Nullstellensatz implies
\((x,y,z)^N\subset(g_1,g_2,g_3)\) for some \(N\). Surjectivity of
\(S_{N-3}^3\to S_N\), \((h_i)\mapsto\sum h_ig_i\), is an open
matrix-rank condition. Conversely, surjectivity excludes a projective
common zero. The complement of the image is the union of these open
sets, and is therefore open.

The image has dimension 29. For example, \(g_1=x^3-z^3\),
\(g_2=y^3-z^3\) meet in nine distinct points, and one can choose \(g_3\)
to vanish at exactly one of them. At that point the first two equations
have an invertible Jacobian in the two affine point coordinates. Locally
they solve for the point, and the third equation imposes one independent
condition on the coefficient space. Thus the image is an irreducible
hypersurface. Since the coefficient ring is a unique factorization
domain, it has an irreducible defining polynomial \(D\), unique up to a
constant.

For a general fixed pair of cubics, there are nine distinct intersection
points. Here is the degree count. Coprime cubics form a regular sequence
of length two: the first is nonzero, and coprimality makes the second a
nonzerodivisor modulo the first, by unique factorization. The two
multiplication exact sequences give \[
\operatorname{Hilb}_{S/(g_1,g_2)}(t)=\frac{(1-t^3)^2}{(1-t)^3}.
\] Its coefficients are eventually nine. For distinct transverse
projective zeros the eventual coefficient equals their number.
Homogeneous interpolation separates the points. Transversality
identifies the localized ideal with the reduced point ideal at each of
them. The difference between the complete-intersection ideal and the
intersection of these point ideals is consequently supported only at the
origin; the Nullstellensatz makes a power of \((x,y,z)\) annihilate this
finitely generated difference, so its graded pieces eventually vanish.
Transversality is an open condition, and the displayed pair exhibits it.
For this general pair, \(D\), viewed as a polynomial in \(g_3\),
vanishes on the nine distinct evaluation hyperplanes \(g_3(P)=0\). Each
occurs with multiplicity one: at a triple having just one common point,
the preceding implicit-function calculation makes that evaluation the
local defining equation, with nonzero derivative in a coefficient of
\(g_3\). We may choose a smooth point of the irreducible hypersurface
here: in characteristic zero an irreducible polynomial cannot divide all
its nonzero partial derivatives, so its smooth locus is dense. Hence
\(D\) has degree nine in that group of coefficients. Symmetry gives
degree nine in each group.

The promised determinant description supplies a rational formula and its
coefficient field. Consider
\protect\phantomsection\label{equation-B30}{}

\[
S_1^3\xrightarrow{K}S_4^3\xrightarrow{M}S_7,
\qquad M(h_1,h_2,h_3)=\sum h_ig_i,
\tag{B30}
\] where the three components of \(K\) send a linear form \(\ell\), for
a pair \(i<j\), to the tuple with \(g_j\ell\) in slot \(i\),
\(-g_i\ell\) in slot \(j\), and zero elsewhere. Its dimensions are
\(9,45,36\), and \(MK=0\). At \((x^3,y^3,z^3)\), \(M\) is surjective,
because a monomial of degree seven must be divisible by one of
\(x^3,y^3,z^3\). Also \(K\) is injective. Thus the generic complex is
exact.

The elementary complementary-minor identity for an exact complex says
\protect\phantomsection\label{equation-B31}{}

\[
\det M_J=\varepsilon_J T\det K_{J^c}
\tag{B31}
\] for every set \(J\) of 36 domain basis elements, where
\(\varepsilon_J\) is the basis-shuffle sign and \(T\) is a common
rational function. To prove the identity, choose a basis of the kernel
of \(M\) from the columns of \(K\), complete it to a basis of the middle
space, and take the top exterior power. Expanding that wedge in the
original basis gives exactly the complementary minors and their shuffle
signs.

In fact \(T\) is polynomial. The nine-row minors of \(K\) have no common
irreducible factor. Whenever \(g_1,g_2\) are coprime, \(K\) is
injective: its third-slot equation is a syzygy between \(g_1,g_2\) with
linear coefficients; divisibility forces both coefficients to be zero
because the cubics have degree three, and the remaining coefficient is
then zero. The locus of noncoprime pairs has codimension at least two.
For a common factor of degree one, two or three, its dimension in the
20-dimensional pair space is at most \[
2+6+6=14,\quad 5+3+3=11,\quad 9+1+1=11,
\] respectively: choose the common factor up to scalar and the two
remaining factors. Consequently no hypersurface can divide all those
minors. If the rational function \(T\) had an irreducible denominator
after cancellation, (B31) would force that denominator to divide every
minor of \(K\), a contradiction.

Choose \(J\) as follows: for every degree-seven monomial, take the
multiple of the first \(g_i\) for which its exponent in variable \(i\)
is at least three. This selects 15, 12 and 9 domain elements in the
three slots. In the complementary minor of \(K\), the remaining three
rows in slot two involve \(g_1\), and the remaining six rows in slot
three give three further \(g_1\) factors and three \(g_2\) factors.
Equivalently its block form is triangular by the pairs
\((1,2),(1,3),(2,3)\). Thus its multidegree is \((6,3,0)\), whereas
\(\det M_J\) has multidegree \((15,12,9)\). It follows that \(T\) has
multidegree \((9,9,9)\).

At a general triple with a common point, evaluation at that point
annihilates the image of \(M\), so all its maximal minors vanish; \(K\)
still has rank nine. Therefore \(T\) vanishes on the irreducible
incidence hypersurface, and \(D\) divides \(T\). Their equal
multidegrees make the quotient a nonzero constant. Normalize the sign
and constant at the three pure cubes. This gives \(\operatorname{Res}\),
proves the formula, and proves the zero-set assertion everywhere,
including multiple intersections. The matrices and normalization are
rational, so the resulting polynomial is rational. The zero-set
assertion extends to every characteristic-zero field as well. The
finitely generated field of the given coefficients embeds in
\(\mathbb C\). For each degree, surjectivity of the multiplication
matrix used above is preserved under extension of fields. The
Nullstellensatz's equivalence between some such surjectivity and absence
of a projective zero therefore transfers the assertion to any algebraic
closure of the coefficient field. \(\square\)

Define \protect\phantomsection\label{equation-B32}{}

\[
I_{27}(F)=\frac{\operatorname{Res}(F_x,F_y,F_z)}{4^{27}}.
\tag{B32}
\] Euler's equation \(4F=xF_x+yF_y+zF_z\) shows that its zero set is
exactly the singular quartics. Its degree is 27, and its value on
\(x^4+y^4+z^4\) is one.

It is an invariant. Changes of variables and invertible combinations of
the three cubic equations preserve the resultant hypersurface, so the
irreducible polynomial in Lemma 5.1 transforms by a character. On either
special linear group that character is trivial. Indeed, its restriction
to an elementary one-parameter subgroup is a regular character
\(\mathbb G_a\to\mathbb G_m\); the only units in \(\mathbb C [t]\) are
nonzero constants, and the identity value is one. Elementary matrices
generate \(\operatorname{SL}_3\). The chain rule for the three partial
derivatives uses these two special linear actions, proving the
assertion.

\subsubsection{The exact elimination
certificates}\label{the-exact-elimination-certificates}

Some coefficients below are large, but the calculation is finite. The
editable source package contains \texttt{quartic\_algebra.py}, the
construction scripts and every polynomial certificate used here. In
particular, the certificates record the following rule for every
contraction: \protect\phantomsection\label{equation-B33}{}

\[
(a_\alpha x^\alpha)\mathbin{\lrcorner}(b_\beta u^\beta)
=a_\alpha b_\beta\frac{\beta!}{(\beta-\alpha)!}u^{\beta-\alpha}
\quad(\beta\geq\alpha),
\tag{B33}
\] with zero contribution otherwise. Equations (B26)--(B29) and (B33)
completely specify the calculation; no source-paper software is called.

We use two elementary certificate principles. If two univariate
polynomials have a common root, their Sylvester determinant is zero.
Evaluation at that root gives a nonzero vector annihilated by the
transpose of the Sylvester map \((a,b)\mapsto ap+bq\); this argument
remains valid if a leading coefficient specializes to zero. Iterating
this observation permits elimination in two stages. If several
univariate polynomials vanish at the same value, their monic greatest
common divisor also vanishes there: the Euclidean algorithm expresses
that greatest common divisor as their polynomial linear combination.

For an ideal-membership certificate, the data consist simply of
polynomials \(h_i\) and the checked equality \(q=\sum h_ip_i\). This
forces \(q=0\) at every common zero of the input polynomials. The stored
branch identities use rational coefficients and are verified by exact
multiplication. They do not require an assumption that numerical sample
points exhaust a variety.

\protect\phantomsection\label{lemma-5-2}{}

\textbf{Lemma 5.2 (the distinct-tangent calculation).} For
\protect\phantomsection\label{equation-B34}{}

\[
F=xyz^2+(x^3+y^3)z+Ay^4+Bxy^3-\frac92x^2y^2+Dx^3y+Ex^4,
\tag{B34}
\] the equations \(I_6=I_9=I_{12}=I_{15}=I_{18}=0\) imply either
\protect\phantomsection\label{equation-B35}{}

\[
A=E=-\frac32,\quad B=D=-\frac{15}{4},
\tag{B35}
\] or \protect\phantomsection\label{equation-B36}{}

\[
4A^2-6A+9=0,\quad E=\frac32-A,\quad
B=\frac{15}{4}-\frac52A,\quad D=\frac52A.
\tag{B36}
\]

\textbf{Proof.} Equation (B29) gives \(I_6=(4AE-9)/82944\), so \(A\ne0\)
and \(E=9/(4A)\). Substitute this value into the four expressions of
degrees 9, 12, 15 and 18. Clear the denominators, obtaining
\(p_0,p_1,p_2,p_3\in\mathbb Q [D,B,A]\), in this order. The coefficient
lists are in \texttt{case-2.1-resultants.json}; they are derived from
(B26)--(B29), not supplied as unexplained inputs.

Put \(r_i=\operatorname{Res}_D(p_0,p_i)\), for \(i=1,2,3\), then form
\(\operatorname{Res}_B(r_i,r_j)\), for the three pairs. Their degrees in
\(A\) are 204, 262 and 298. Exact Euclidean division gives their monic
greatest common divisor \protect\phantomsection\label{equation-B37}{}

\[
\frac{A^{66}(2A+3)^8(4A^2-6A+9)^8}{2^{24}}.
\tag{B37}
\] Every common zero of the original equations makes all three final
resultants vanish, so (B37), together with \(A\ne0\), forces \(2A+3=0\)
or \(4A^2-6A+9=0\).

For the first branch, substitute \(A=E=-3/2\) into the original four
polynomial equations. Exact ideal-membership identities give \[
B+D+\frac{15}{2}=0,\qquad
\left(B+\frac{15}{4}\right)^2=0.
\] These imply (B35). For the second branch, include \(4A^2-6A+9\) among
the input equations and substitute \(E=3/2-A\). The corresponding
identities give \[
\frac23AB-5A-B+D=0,\qquad
\left(B-\frac{15}{4}+\frac52A\right)^2=0.
\] The square determines \(B\); reducing the linear equation modulo
\(4A^2-6A+9\) determines \(D\). This proves (B36). The full multipliers
witnessing both pairs of identities are in
\texttt{branch-identities.json}, and the verification is direct
substitution and multiplication. \(\square\)

Both forms in this lemma are a line times a cubic with a flex on that
line. For (B35), \protect\phantomsection\label{equation-B38}{}

\[
F=-\frac14(3x+3y-2z)
(2x^3+3x^2y+3xy^2+2y^3+2xyz).
\tag{B38}
\] At \(P=(-1,1,0)\) the cubic gradient is \((3,3,-2)\), and restricting
the cubic to \(z=3(x+y)/2\) gives \(2(x+y)^3\). Thus it is smooth at
\(P\) and the line meets it three times there.

For (B36), put \[
\begin{aligned}
L&=x+(2A/3-1)y+4Az/9,\\
Q&=x^3+(A-3/2)x^2y-Axy^2+y^3+xyz.
\end{aligned}
\] Exact multiplication modulo \(4A^2-6A+9\) gives \(F=(3/2-A)LQ\). At
\(P=(1-2A/3,1,0)\), the gradient of \(Q\) is
\((-A,3/2,1-2A/3)=-A\nabla L\), which is nonzero. Substitution of
\(L=0\) into \(Q\) gives \((x+(2A/3-1)y)^3\). This verifies the flex and
all scalar factors in the factorizations.

\subsubsection{Why their common zeros are the null
cone}\label{why-their-common-zeros-are-the-null-cone}

A convenient sufficient condition for a form to lie in the null cone is
the following. If integers \((w_x,w_y,w_z)\) sum to zero and every
monomial in \(F\) has positive weight \(iw_x+jw_y+kw_z\), then
\(F(t^{w_x}x,t^{w_y}y,t^{w_z}z)\to0\) as \(t\to0\). This substitution is
the inverse action of a determinant-one diagonal matrix. Every
positive-degree homogeneous invariant therefore vanishes on \(F\).

\protect\phantomsection\label{lemma-5-3}{}

\textbf{Lemma 5.3 (quartic common zeros).} The common zero set of
\(I_3,I_6,I_9,I_{12},I_{15},I_{18},I_{27}\) is the null cone.

\textbf{Proof.} We work over an algebraically closed characteristic-zero
field. The zero form causes no difficulty. For a nonzero common zero,
(B32) supplies a singular point. General linear coordinate changes and
nonzero scalar multiples preserve the vanishing of homogeneous
invariants: over this field a general linear matrix is a scalar times a
special linear matrix. Move the singular point to \((0,0,1)\). The
coefficients of \(z^4,xz^3,yz^3\) vanish, so
\protect\phantomsection\label{equation-B39}{}

\[
F=q_2(x,y)z^2+q_3(x,y)z+q_4(x,y).
\tag{B39}
\] If \(q_2=0\), the point has multiplicity at least three; weights
\((1,1,-2)\) are positive on every monomial in (B39). Hence this form is
in the null cone.

Suppose the tangent quadratic has rank one. Normalize it to \(x^2\). If
the coefficient of \(y^3\) in \(q_3\) is zero, the form has the shape
\protect\phantomsection\label{equation-B40}{}

\[
F=x^2z^2+(Lxy^2+Mx^2y+Nx^3)z
+Ay^4+Bxy^3+Cx^2y^2+Dx^3y+Ex^4.
\tag{B40}
\] Direct use of (B26)--(B29) gives \[
I_3=(6A+L^2)/72,\qquad
I_6=L^2(24A-L^2)/746496.
\] Thus \(A=-L^2/6\), and \(I_6=-5L^4/746496\) forces \(A=L=0\). Weights
\((4,-1,-3)\) are now positive on all the remaining monomials in (B40).

If the coefficient of \(y^3\) is nonzero, scale it to one. Replacing
\(y\) by \(y-Lx/3\) removes \(xy^2z\); subsequently replacing \(z\) by
\(z+bx+cy\), with suitable \(b,c\), removes \(x^3z,x^2yz\). The
resulting shape is \protect\phantomsection\label{equation-B41}{}

\[
F=x^2z^2+y^3z+Ay^4+Bxy^3+Cx^2y^2+Dx^3y+Ex^4.
\tag{B41}
\] Here \(I_3=A/12\) and \(I_6=C/20736\), so \(A=C=0\). With these
substitutions, \[
I_{15}=\frac{1125D^2}{165112971264},\qquad
I_9=\frac{801B^2+585E}{2985984}.
\] Thus \(D=0\), \(E=-89B^2/65\). Substitution into \(I_{18}\) gives \[
I_{18}=\frac{6251B^4}{114831654912},
\] forcing \(B=E=0\). The remaining form \(x^2z^2+y^3z\) has positive
weights \((-3,-1,4)\).

It remains to treat rank two. Normalize \(q_2=xy\), and shift \(z\) by a
linear form to remove the terms \(xy^2z,x^2yz\). Write
\protect\phantomsection\label{equation-B42}{}

\[
F=xyz^2+(sy^3+tx^3)z+Ay^4+Bxy^3+Cx^2y^2+Dx^3y+Ex^4.
\tag{B42}
\] If \(s,t\ne0\), diagonal scaling makes both one. Explicitly, choose
\(\alpha^3/\beta^3=s/t\) and \(\gamma=s\beta^2/\alpha\); substitution by
\((\alpha x,\beta y,\gamma z)\), followed by division by
\(\alpha\beta\gamma^2\), does so. The equation \(I_3=0\) gives
\(C=-9/2\), and Lemma 5.2 applies. Its factorizations show a line \(L\)
times a cubic smooth at a point where \(L\) has triple contact. Send
that point to \((0,0,1)\) and the line to \(x=0\). Smoothness and triple
contact put the product in (B40) with \(A=L=0\): the cubic has a nonzero
\(xz^2\) coefficient, no \(yz^2,y^2z\) terms, and restricts to a
multiple of \(y^3\) on \(x=0\). The same weights \((4,-1,-3)\) prove
null-cone membership.

If precisely one of \(s,t\) is nonzero, exchange \(x,y\) if necessary
and scale to \(s=0,t=1\). The first two parameter equations give \[
C=0,\qquad AE=0.
\] If \(A=0\), \(I_9=0\) forces \(B^2E=0\). Assuming \(B\ne0\) would
give \(E=0\), and \(I_{12}=0\) then gives \(D^2=3B\). The remaining
degree-18 equation is a nonzero scalar times \(B^3D(9B+D^2)=0\); hence
\(D=0\), contradicting \(B\ne0\). Therefore \(B=0\).

If instead \(A\ne0\), then \(E=0\), and \(I_9=I_{15}=0\) give \[
B=\frac{2D^2}{45},\qquad A^2=\frac{42}{125}B^2D.
\] After these substitutions the numerator of \(-34828517376I_{12}\) is
\(5216D^6/375\). Thus \(D=B=0\), then \(A=0\), a contradiction. We have
proved \(A=B=C=0\) in this case. The form has multiplicity at least
three at \((0,1,0)\), and weights \((1,-2,1)\) give the required limit.

Finally, if \(s=t=0\), the equations of degrees 3, 6 and 12 give \[
C=0,\qquad AE=0,\qquad BD=0;
\] the degree-nine equation gives \(AD^2+B^2E=0\). If \(A\ne0\), then
\(E=D=0\). If \(A=0\), either \(B=0\), or \(B\ne0\) forces \(E=D=0\).
Consequently either \(A=B=0\) or \(E=D=0\). Exchanging \(x,y\) reduces
the second alternative to the first, and weights \((1,-2,1)\) again
apply.

These exhaust (B39). Thus every common zero is in the null cone.
Conversely, all seven expressions are homogeneous positive-degree
invariants, so they vanish on the null cone by its definition.
\(\square\)

\subsubsection{Independence and module
finiteness}\label{independence-and-module-finiteness}

\protect\phantomsection\label{theorem-5-4}{}

\textbf{Theorem 5.4 (Dixmier's parameters).} Over every
characteristic-zero field the seven invariants (B28), (B29) and (B32)
are a homogeneous system of parameters for ternary quartics. Their
degrees are \(3,6,9,12,15,18,27\).

\textbf{Proof.} Lemma 5.3 and
\href{NOE-INV-01.html\#theorem-11-1}{Hilbert's nullcone finiteness
theorem, Theorem 11.1 in the first lesson}, prove module finiteness over
\(\mathbb C\). It remains to prove algebraic independence; the following
small exact witness does so without a quotient-dimension assumption.

Order the ordinary coefficients by the monomials \[
x^4,x^3y,x^3z,x^2y^2,x^2yz,x^2z^2,xy^3,
xy^2z,xyz^2,xz^3,y^4,y^3z,y^2z^2,yz^3,z^4.
\] Evaluate at the integer coefficient vector
\protect\phantomsection\label{equation-B43}{}

\[
(-4,-5,9,3,4,3,-7,7,2,-1,7,-1,-9,6,-3).
\tag{B43}
\] Take the first seven columns of the seven-parameter Jacobian in
parameter order \(3,6,9,12,15,18,27\). All denominators of its rational
entries are prime to 101. Reduction modulo 101 gives
\protect\phantomsection\label{equation-B44}{}

\[
\begin{pmatrix}
63&55&90&49&92&14&37\\
26&43&98&63&99&97&90\\
90&62&48&57&87&86&36\\
5&98&56&87&24&69&97\\
17&58&82&52&33&21&60\\
85&70&97&88&4&96&56\\
54&76&72&40&70&79&23
\end{pmatrix},
\qquad\det\equiv47\pmod{101}.
\tag{B44}
\] The first six rows follow directly from the contractions. To specify
the last row completely, form the matrices in (B30), selecting the 36
domain monomials by the rule in Lemma 5.1 and taking the complementary
nine rows of \(K\). Both selected determinants are nonzero at (B43). If
their matrices are \(M_0,K_0\), the derivative of their quotient in
coefficient direction \(h\) is
\protect\phantomsection\label{equation-B45}{}

\[
I_{27}(F)\left(\operatorname{tr}(M_0^{-1}M_h)
-\operatorname{tr}(K_0^{-1}K_h)\right).
\tag{B45}
\] This is the usual determinant derivative, obtained by expanding
\(\det(I+tA)=1+t\operatorname{tr}A+O(t^2)\). Normalize the quotient by
its value at the Fermat quartic, as in (B32). The certificate
\texttt{parameter-independence.json} contains the integer matrices, all
rational derivatives and the minor;
\texttt{check\_quartic\_independence.py} constructs them from
(B25)--(B33). Thus (B44) is an exact rational nonvanishing certificate,
although its compact display uses a finite field.

A nonzero Jacobian minor proves independence in characteristic zero. For
completeness, choose a nonzero polynomial relation of least total
degree, if one exists. Differentiate it after substitution of the seven
parameters. Full Jacobian rank forces all its partial derivatives, after
that same substitution, to vanish. Some partial derivative is a nonzero
polynomial of smaller degree, contradicting the choice of the relation.

Finally, descend finiteness from \(\mathbb C\) to \(\mathbb Q\) and
extend to \(k\). In each coefficient degree, invariants are the
simultaneous kernels of the six rational elementary-matrix derivations
for \(\operatorname{SL}_3\); the unipotent argument of
\href{NOE-INV-01.html\#theorem-2-1}{Theorem 2.1}, with Gaussian
elimination in dimension three, proves this criterion. These finite
kernels commute with field extension. The quotient by the ideal
generated by the seven parameters also commutes with field extension
degree by degree. Its complex dimension is finite, so its rational
dimension is finite. Choose homogeneous rational lifts of a basis of
that quotient; the degree induction in
\href{NOE-INV-01.html\#theorem-11-1}{Theorem 11.1 of the first lesson}
makes them module generators. Extending scalars gives finiteness over
\(k\), while (B44), as a rational nonzero determinant, gives
independence over every characteristic-zero field. \(\square\)

The parameter theorem is complete. The next section proves the full
generator theorem by identifying an explicit free module basis.

\subsection{6. Thirteen generators for a ternary
quartic}\label{section-6}

The parameter theorem leaves a concrete question: which elements
generate its finite module? We answer it by constructing fifty
independent elements with exactly the Hilbert-series degrees. This also
proves algebra generation. The calculation below is an exact finite
certificate, followed by an argument that applies in every degree and
over every characteristic-zero field.

Dixmier supplied the seven parameters; Ohno supplied the additional
generators and the fifty monomials used below. Compare
\href{https://aeb.win.tue.nl/math/ohno-preprint.2007.05.15.pdf}{Ohno's
primary preprint, Definition 3.1 and Theorem 4.1}, and
\href{https://math.uni-paderborn.de/fileadmin-eim/mathematik/AG-Computeralgebra/Publications-elsenhans/inv_pl_quart_5a.pdf}{Elsenhans,
Sections 5--6}. The polynomial formulas and arithmetic certificates
below give the proof.

\subsubsection{Six additional
contractions}\label{six-additional-contractions}

Retain the conventions and contractions (B25)--(B33). Define
\protect\phantomsection\label{equation-B46}{}

\[
\begin{aligned}
H&=\frac{\det(F_{x_i x_j})}{1728},&
\xi&=\frac{\sigma\mathbin{\lrcorner}H}{72},&
\pi&=\frac{\rho\mathbin{\lrcorner}H}{2},\\
\eta&=\frac{\xi\mathbin{\lrcorner}\sigma}{12},&
\nu&=\frac{\eta\mathbin{\lrcorner}\pi}{4}.
\end{aligned}
\tag{B46}
\] The Hessian \(H\) is a degree-three covariant of order six. Indeed,
the chain rule transforms its Hessian matrix by congruence, and the two
determinant factors equal one on the special linear group. The forms
\(\xi,\pi,\eta,\nu\) have coefficient degrees \(5,7,7,14\) and orders
\(2,4,2,2\), respectively. The first and last are covariants, \(\eta\)
is a contravariant, and \(\pi\) is a covariant. These types follow by
alternating the dual spaces in each equivariant contraction.

Write their quadratic matrices as \(C,E,N\), so \(\xi=x^{\mathsf t}Cx\),
\(\eta=u^{\mathsf t}Eu\), \(\nu=x^{\mathsf t}Nx\). With \(A,B\) still
the matrices of \(\rho,\tau\), put
\protect\phantomsection\label{equation-B47}{}

\[
\begin{aligned}
J_9&=\operatorname{tr}(CA),&J_{12}&=\operatorname{tr}(BE),\\
J_{15}&=\det C,&J_{18}&=\operatorname{tr}(\operatorname{adj}C\operatorname{adj}A),\\
J_{21}^{a}&=\det E,&J_{21}^{b}&=\operatorname{tr}(NE).
\end{aligned}
\tag{B47}
\] The congruence and inverse-congruence calculation following (B29)
proves that all six are invariants. Their subscripts give their
coefficient degrees. In particular, (B46)--(B47) are definitions over
\(\mathbb Q\), with no choice of roots or auxiliary normal form. The
script \texttt{quartic\_generators.py} implements precisely these
contractions.

\subsubsection{The Hilbert series in all
degrees}\label{the-hilbert-series-in-all-degrees}

\protect\phantomsection\label{lemma-6-1}{}

\textbf{Lemma 6.1 (the quartic character calculation).} For
\(R=\mathbb C [c_{ijk}]^{\operatorname{SL}_3}\),
\protect\phantomsection\label{equation-B48}{}

\[
\operatorname{Hilb}_R(t)=
\frac{N(t)}{(1-t^3)(1-t^6)(1-t^9)(1-t^{12})(1-t^{15})(1-t^{18})(1-t^{27})},
\tag{B48}
\] where \protect\phantomsection\label{equation-B49}{}

\[
\begin{aligned}
N(t)={}&1+t^9+t^{12}+t^{15}+2t^{18}+3t^{21}+2t^{24}+3t^{27}\\
&+4t^{30}+3t^{33}+4t^{36}+4t^{39}+3t^{42}+4t^{45}\\
&+3t^{48}+2t^{51}+3t^{54}+2t^{57}+t^{60}+t^{63}+t^{66}+t^{75}.
\end{aligned}
\tag{B49}
\]

\textbf{Proof.} Averaging over \(\operatorname{SU}(3)\) projects each
finite-dimensional coefficient-degree space onto its special-linear
invariants, by the density and Reynolds arguments in
\href{NOE-INV-01.html\#section-8}{the first lesson, Sections 8--9}. The
trace of that projection is its rank. A torus element
\(\operatorname{diag}(x,y,(xy)^{-1})\) has the fifteen weights \[
\mathcal W=\{(i-k,j-k):i,j,k\geq0,\ i+j+k=4\}.
\] The symmetric-algebra character is therefore
\(\prod_{(a,b)\in\mathcal W}(1-tx^ay^b)^{-1}\), as follows by
multiplying the geometric series of its fifteen generators. For
\(|t|<1\) this series converges uniformly on the unit torus, so its
averaging can be performed term by term.

We justify the torus integration factor here. Every unitary matrix is
unitarily diagonalizable. For matrices with three distinct eigenvalues,
the map from a conjugating flag and a diagonal matrix to its conjugate
is six to one, the six choices being the eigenvalue order. In the
tangent plane corresponding to a pair of eigenvectors \(i,j\), its
differential multiplies the two-dimensional real area by
\(|z_i-z_j|^2\): in complex coordinates it is multiplication by
\(z_j-z_i\). Thus its Jacobian is \(\prod_{i<j}|z_i-z_j|^2\). Matrices
with a repeated eigenvalue form a lower-dimensional set and have measure
zero; this follows also from the nonzero discriminant equation in real
coordinate charts. Change of variables for the invariant volume
consequently expresses a class-function integral as a constant times its
torus integral with this Jacobian. The constant is determined by the
integral of one. Expanding the determinant \(\det(z_i^{j-1})\) and its
conjugate shows that the constant term of their squared modulus is six:
the only matching exponent permutations are equal permutations, modulo
the common determinant-one exponent shift. Hence the normalized factor
is \(1/6\).

Put \[
D_+=(1-x^{-1}y)(1-x^{-2}y^{-1})(1-x^{-1}y^{-2}).
\] The squared Vandermonde on the unit torus is the product of these
factors and their inverses in \(x,y\). Direct expansion gives
\(\sum_{w\in S_3}wD_+=\prod_{i\ne j}(1-z_i/z_j)\). For a symmetric torus
character, all six constant terms in that sum are equal. Therefore
normalized averaging is the single constant term
\protect\phantomsection\label{equation-B50}{}

\[
\operatorname{Hilb}_R(t)=\operatorname{CT}_{x,y}
\frac{D_+}{\prod_{(a,b)\in\mathcal W}(1-tx^ay^b)}.
\tag{B50}
\] This establishes the counting formula without assuming a freeness or
Cohen--Macaulay theorem.

Here is a finite rational evaluation of (B50), specifying the
calculation rather than extrapolating a list of low degrees. First hold
\(|y|=1\) and take \(|t|\) sufficiently small. The poles inside the
\(x\)-circle come from the six weights \[
(-1,-1),\ (-2,-3),\ (-1,2),\ (-2,0),\ (-3,-2),\ (-4,-4).
\] There is no pole at \(x=0\): clearing the six negative-weight
denominators gives total \(x\)-exponent thirteen, whereas \(D_+\) has
negative exponent at most four; after the additional \(1/x\) the
exponent is still at least eight. A weight \((-r,b)\) gives roots
\(\alpha^r=ty^b\). Its residue in the integrand divided by \(x\) is
\protect\phantomsection\label{equation-B51}{}

\[
\frac1r\frac{D_+(\alpha,y)}{
\prod_{(a,c)\ne(-r,b)}(1-t\alpha^ay^c)}.
\tag{B51}
\] Sum these \(r\) residues by reducing the fraction in the algebra
\(\mathbb Q(t,y)[\alpha]/(\alpha^r-ty^b)\). The normalized trace is its
coefficient of \(1\): the other basis monomials have trace zero. To
compute that coefficient, form the \(r\)-by-\(r\) matrix for
multiplication by the denominator, and solve against the numerator by
Cramer's rule. All \(r\) are at most four. Negative powers are reduced
using \(\alpha^r=ty^b\); multiplication by a common monomial clears
them. This is the complete algorithm in \texttt{quartic\_hilbert.py}.

Adding the six rational functions and cancelling their polynomial
greatest common divisors gives a numerator with 294 terms and a
denominator with 726 terms in \(\mathbb Q [t,y]\). Their full
coefficients are in \texttt{hilbert-series.json}. Factoring the latter
denominator, its factors whose roots tend to zero with \(t\) are
\protect\phantomsection\label{equation-B52}{}

\[
\begin{gathered}
t-y,\quad t+y^2,\quad t^2-y,\quad t^2+ty+y^2,\\
t^2+ty^2+y^4,\quad t^5-y^4,\quad t^5-y^7,\quad t^7-y^8,
\end{gathered}
\tag{B52}
\] each once, and \(t-y^2\) twice. All other nonconstant \(y\)-factors
have roots tending to infinity; factors involving only \(t\) create no
\(y\)-pole. These assertions follow by substituting \(y=t^{a/b}v\) in
each binomial or homogeneous cyclotomic factor; they also determine the
inside roots without numerical root finding.

There is now also the residue at \(y=0\). If the first rational function
is \(n(y)/d(y)\), it is \(n(0)/d(0)\). For any monic inside factor \(q\)
of multiplicity \(m\), write \(d=q^m h\), and put \(\phi=n/(yh)\). For
\(m=1\), the sum of the residues is the trace of \(\phi/q'\) in
\(\mathbb Q(t)[y]/q\). For \(m=2\), it is the trace of
\protect\phantomsection\label{equation-B53}{}

\[
\frac{\phi'}{(q')^2}-\frac{\phi q''}{(q')^3}.
\tag{B53}
\] Indeed, at a simple root \(\alpha\), write \(q=(y-\alpha)q_1\); then
\(q_1(\alpha)=q'(\alpha)\) and \(q_1'(\alpha)=q''(\alpha)/2\). Taking
the derivative of \(\phi/q_1^2\) proves (B53). The derivative \(\phi'\)
is computed before reducing modulo \(q\). Inverses are obtained by the
polynomial Euclidean algorithm, and a trace is the sum of the diagonal
entries of the multiplication matrix. These quotient degrees are at most
eight.

The script \texttt{quartic\_hilbert\_second.py} performs those nine
traces and the zero residue with rational polynomial pairs. Its
certificate contains every traced fraction. Adding them gives
\(n_*(t)/d_*(t)\); exact multiplication verifies
\protect\phantomsection\label{equation-B54}{}

\[
n_*(t)\prod_{d\in\{3,6,9,12,15,18,27\}}(1-t^d)
=d_*(t)N(t).
\tag{B54}
\] Thus the rational functions are equal. Equality for sufficiently
small complex \(t\) gives equality of their Taylor series and proves
(B48) in all degrees. This is the classical Hilbert series credited to
Shioda in Ohno's Theorem 1.1; (B50)--(B54) give its full independent
verification here. \(\square\)

\subsubsection{Fifty monomials with the required
degrees}\label{fifty-monomials-with-the-required-degrees}

For this table only, abbreviate the six additional invariants by \[
(u,v,w,r,s,h)=(J_9,J_{12},J_{15},J_{18},J_{21}^{a},J_{21}^{b}).
\] Each entry in the right column is a separate candidate secondary
\(W\). We order by increasing degree, with the displayed left-to-right
order within each row. There are fifty entries, including one.

{\def\LTcaptype{none} % do not increment counter
\begin{longtable}[]{@{}ll@{}}
\toprule\noalign{}
Degree & Secondary monomials \\
\midrule\noalign{}
\endhead
\bottomrule\noalign{}
\endlastfoot
0 & \(1\) \\
9 & \(u\) \\
12 & \(v\) \\
15 & \(w\) \\
18 & \(u^2,r\) \\
21 & \(uv,s,h\) \\
24 & \(uw,v^2\) \\
27 & \(u^3,ur,vw\) \\
30 & \(u^2v,us,uh,vr\) \\
33 & \(u^2w,uv^2,wr\) \\
36 & \(u^2r,uvw,ws,wh\) \\
39 & \(u^3v,u^2s,u^2h,rh\) \\
42 & \(uwr,v^2r,sh\) \\
45 & \(u^2vw,uws,uwh,vwr\) \\
48 & \(u^3s,u^3h,urh\) \\
51 & \(u^2wr,ush\) \\
54 & \(u^2ws,u^2wh,wrh\) \\
57 & \(u^2rh,wsh\) \\
60 & \(u^2sh\) \\
63 & \(uwrh\) \\
66 & \(uwsh\) \\
75 & \(u^2wsh\) \\
\end{longtable}
}

Counting the entries of each degree gives exactly (B49). The assertion
that they are independent over the parameter ring is the substantive
issue, proved next.

\subsubsection{Why a nonsingular finite-field fibre proves
independence}\label{why-a-nonsingular-finite-field-fibre-proves-independence}

\protect\phantomsection\label{lemma-6-2}{}

\textbf{Lemma 6.2 (lifting an independence certificate).} Let
\(q_1,\ldots,q_m\) and \(W_1,\ldots,W_\ell\) be rational polynomials in
\(n\) variables whose coefficient denominators are prime to a prime
\(p\). Suppose a finite reduced algebra over \(\mathbb F_p\) supplies
points all having the same \(q\)-tuple \(c\in\mathbb F_p^m\). Suppose
the \(q\)-Jacobian has rank \(m\) at every point, and the \(\ell\)
functions \(W_i\) are linearly independent in that finite algebra. Then
the \(W_i\) are linearly independent over
\(\mathbb Q [q_1,\ldots,q_m]\).

\textbf{Proof.} A squarefree polynomial quotient splits into a product
of fields over a finite extension, by the Euclidean algorithm and the
Chinese remainder identity. More generally the assumed reduced finite
algebra splits into a product of copies of a finite extension: each
finite field extension is separable because finite fields are perfect,
as the Frobenius map is bijective. After adjoining all point
coordinates, the evaluations of the \(W_i\) still have column rank
\(\ell\). Select \(\ell\) evaluation rows with a nonzero square
determinant.

We describe the elementary characteristic-zero lifting step. Let
\(\mathbb Z_p=\varprojlim_a\mathbb Z/p^a\mathbb Z\). Successive residues
give a unique base-\(p\) expansion; its first nonzero digit shows that
this is a domain of characteristic zero, with units precisely the
elements nonzero modulo \(p\). For a finite residue field
\(\mathbb F_{p^f}\), choose its monic irreducible defining polynomial
and any monic integer lift. The algebra obtained by adjoining a root
over \(\mathbb Z_p\) is complete, free of rank \(f\), and has that
residue field. Every element not divisible by \(p\) is a unit: lift its
residue inverse and invert \(1-pa\) by the convergent geometric series.
Factoring out the first nonzero \(p\)-power from each element shows that
products of nonzero elements are nonzero. Denote this complete
characteristic-zero domain by \(O\).

At each selected point choose \(m\) variable directions with invertible
Jacobian minor, and lift the other coordinates arbitrarily to \(O\).
Choose the same integer lift \(\widetilde c\) of the primary tuple for
every point. If a current coordinate tuple solves \(q=\widetilde c\)
modulo \(p^a\), a correction \(p^a v\) solves it modulo \(p^{a+1}\) by
the invertible Jacobian and the first-order polynomial expansion. Start
modulo \(p\) and repeat. Completeness gives a tuple with
\(q=\widetilde c\) exactly. The chosen minor remains a unit.

Now introduce independent formal variables \(z_1,\ldots,z_m\) and solve
\(q=\widetilde c+z\) over \(O [[z_1,\ldots,z_m]]\) at each of these
lifted points. The same linear correction, this time by total formal
degree, constructs the solution recursively: after degrees below \(a\)
have been solved, the degree-\(a\) error is removed by the invertible
constant Jacobian. The square evaluation matrix of \(W\) has unit
constant determinant, since its reduction modulo \(p\) is the selected
nonzero determinant. It is therefore invertible over this formal
power-series ring.

If \(\sum f_i(q)W_i=0\) with \(f_i\in\mathbb Q [T_1,\ldots,T_m]\), clear
their common integer denominator. Evaluation at all these formal tuples
and inversion of the matrix gives \(f_i(\widetilde c+z)=0\). Translation
by a constant is injective on polynomial rings over the
characteristic-zero domain \(O\); hence all \(f_i\) vanish. This proves
independence. \(\square\)

The full-rank Jacobian hypothesis is essential. A rank calculation in an
arbitrary specialized fibre alone need not justify a characteristic-zero
conclusion.

\subsubsection{The exact fibre and its rank
certificate}\label{the-exact-fibre-and-its-rank-certificate}

Work modulo \(101\), where all contraction denominators are units. Use
the family \protect\phantomsection\label{equation-B55}{}

\[
F_{A,B,D}=xyz^2+(x^3+y^3)z+Ay^4+Bxy^3-\frac92x^2y^2+Dx^3y+\frac9{4A}x^4.
\tag{B55}
\] Here \(A\ne0\); the contraction table for (B34) gives \(I_3=I_6=0\).
Let \(q_d=I_d(F_{A,B,D})\). Take the numerator \(p_0\) of \(q_9\) and
the numerator \(p_1\) of \(q_{18}\), in reduced rational form with
integer coefficients. Require
\protect\phantomsection\label{equation-B56}{}

\[
p_0=p_1=0,\qquad q_{12}^5-q_{15}^4=0,\qquad q_{12}q_{15}\ne0.
\tag{B56}
\] The complete reduced numerators and denominators are supplied by
\texttt{case-2.1.json} and the substitution \(E=9/(4A)\), so
(B28)--(B29) also define them without relying on that table. Put
\(a=q_{12}/q_{15}\). Then (B56) gives \(a^4q_{12}=a^5q_{15}=1\). Choose
\(\lambda^3=a\) over a finite extension and scale \(F\) by \(\lambda\).
Its seven parameter values become
\protect\phantomsection\label{equation-B57}{}

\[
(I_3,I_6,I_9,I_{12},I_{15},I_{18},I_{27})=(0,0,0,1,1,0,0).
\tag{B57}
\] For a secondary of degree \(d\), its scaled value is \(a^{d/3}W(F)\);
every displayed secondary degree is divisible by three, so no root
choice enters these values.

Here are the precise finite data and checks. The generator certificate
manifest links the full polynomial arrays and their hashes.
\texttt{fibre-secondary-rank.json} supplies a monic polynomial \(G(A)\)
of degree 300, together with coefficient arrays for \(B=b(A)\) and
\(D=d(A)\) in \protect\phantomsection\label{equation-B58}{}

\[
\mathcal A=\mathbb F_{101}[A]/(G).
\tag{B58}
\] Array index \(i\) is the coefficient of \(A^i\); every coefficient is
an integer between zero and 100. Direct polynomial remainder
calculations verify (B56). They also verify that \(G'\), \(A\),
\(q_{12}\), \(q_{15}\), and the following two polynomials are units in
\(\mathcal A\): \protect\phantomsection\label{equation-B59}{}

\[
\begin{aligned}
J&=\det\frac{\partial(p_0,p_1,q_{12}^5-q_{15}^4)}{\partial(A,B,D)},\\
\Delta_h&=\operatorname{Res}_U(h,h'),\\
h(U)&=AU^6-9U^5-4ADU^4+20AU^3-4ABU^2-4A^2U+A.
\end{aligned}
\tag{B59}
\] For every claimed unit the certificate includes its inverse array;
multiplication and reduction give one. It additionally checks the unit
coefficient of the linear \(D\)-remainder used in constructing the
fibre. The inverse of \(G'\) proves that \(G\) is squarefree, so (B58)
is reduced.

The construction is short enough to reproduce. Eliminate \(D\) from
\(p_0,p_1\) by a Sylvester resultant. After discarding the factor
\(A^2\), its degree in \(B\) is twelve. Division of \(p_1\) by \(p_0\)
gives a remainder linear in \(D\), so \(D=-M(A,B)/L(A,B)\) in
\(\mathbb F_{101}(A)[B]/R\). Reduce \(q_{12}^5-q_{15}^4\) in this
degree-twelve algebra and take the determinant of its multiplication
matrix. The resulting rational norm has numerator degree 396 and
denominator degree 198 in \(A\). A Euclidean remainder of degree one in
\(B\) gives its rational coordinate. Taking the squarefree numerator,
removing denominator factors and removing zeros of the two normalizing
invariants leaves exactly degree 300. Original-equation remainders,
\(J\), and \(\Delta_h\) then pass directly. \texttt{quartic\_fibre.py}
implements this construction; the verifier does not assume its
elimination steps are correct, because it recomputes all retained
equations and units directly in (B58).

We explain why these unit checks supply the Jacobian hypothesis of Lemma
6.2. First allow the coefficient \(C\) of \(x^2y^2\), the coefficient
\(E\) of \(x^4\), and the scalar \(\lambda\) to vary. In the
distinct-tangent family, \(I_3=(2C+9)/144\); at \(C=-9/2\),
\(I_6=(4AE-9)/82944\). Thus \(I_3,I_6\) successively solve for \(C,E\),
with derivatives \(1/72\) and \(A/20736\), both units. On that slice,
\(p_0,p_1\) are unit multiples of \(q_9,q_{18}\) at their zeros. The
unit \(J\) says that their differentials together with
\(5q_{12}^4dq_{12}-4q_{15}^3dq_{15}\) have rank three in \(A,B,D\).
Scalar variation adds the vector proportional to \((12q_{12},15q_{15})\)
in the \(q_{12},q_{15}\) coordinates. It is nonzero, and the last
displayed differential annihilates it by (B56). Consequently the four
functions \(I_9,I_{18},I_{12},I_{15}\) have rank four after the scalar
direction is added. Together with \(I_3,I_6\), their rank is six.

Finally, (B55) has an ordinary node at \((0:0:1)\), since its affine
quadratic term is \(xy\). There is no other singularity on \(xy=0\),
because there \(F_z=x^3+y^3\) is nonzero away from that point. On the
chart \(y=1,x=U\ne0\), completing the square in \(z\) gives the binary
polynomial \(h(U)/A\). A further singularity is equivalent to a repeated
root of \(h\). Hence the unit \(\Delta_h\) proves that this is the
unique singular point.

At a quartic with a unique ordinary node, the discriminant has nonzero
differential in the \(z^4\)-coefficient direction. Locally solve the two
affine critical-point equations by the invertible \(xy\)-Hessian.
Euler's equation makes the remaining derivative equation four times the
value of \(F\) at that critical point. Changing the \(z^4\) coefficient
changes that value to first order by one. In the incidence construction
of Lemma 5.1, the unique common point and its two transverse equations
give a smooth local hypersurface with precisely this equation. The
resultant is its reduced defining polynomial, so this local factor has
multiplicity one.

For clarity, the same resultant argument applies over
\(\overline{\mathbb F}_{101}\). The incidence bundle is irreducible of
dimension 29; the pair \(x^3-z^3,y^3-z^3\) still has nine distinct
transverse points, so the multidegree and reduced simple-zero
calculation are unchanged. The determinant quotient in (B31) is integral
before normalization: in the unique factorization ring of integer
coefficient polynomials, a denominator would divide every maximal minor
of \(K\). No nonconstant polynomial can do so by the coprime-pair
argument, and no integer prime can do so because at the three pure cubes
a complementary minor is one. Thus the normalized quotient reduces to
the same resultant modulo 101. Only the normalization \(4^{27}\) is
inverted. This proves the asserted nonzero discriminant derivative also
in the field used by the certificate.

The six family directions preserve the fixed node and thus annihilate
\(dI_{27}\). The extra \(z^4\) direction supplies the seventh rank. The
full seven-parameter Jacobian therefore has rank seven at every point of
the reduced fibre (B58), including after scaling by \(\lambda\).

Evaluate the fifty scaled secondaries in \(\mathcal A\), and write their
coefficient vectors as columns of a \(300\)-by-\(50\) matrix \(V\). The
complete columns are in \texttt{fibre-secondary-rank.json}. Its rows
numbered \[
0,3,6,\ldots,147
\] form a square minor, with
\protect\phantomsection\label{equation-B60}{}

\[
\det V_{\{0,3,\ldots,147\},\{1,\ldots,50\}}=70\quad\text{in }\mathbb F_{101}.
\tag{B60}
\] These are zero-based polynomial coefficient indices; the columns use
the monomial table's order. \texttt{verify\_quartic\_generators.py}
reconstructs (B55), all twelve contractions, all original equations and
unit witnesses, the fifty columns, and (B60). Gaussian elimination
modulo 101 uses only row operations and inverses of nonzero field
elements. The nonzero determinant proves independence in the reduced
algebra; it is an exact calculation, not a probabilistic sample or a
numerical tolerance test. Lemma 6.2 and the full Jacobian just proved
give independence of all fifty monomials over the rational parameter
ring.

\subsubsection{Completion of the generator
theorem}\label{completion-of-the-generator-theorem}

\protect\phantomsection\label{theorem-6-3}{}

\textbf{Theorem 6.3 (Dixmier--Ohno generation).} Over every field \(k\)
of characteristic zero, \protect\phantomsection\label{equation-B61}{}

\[
k [c_{ijk}]^{\operatorname{SL}_3}=
k [I_3,I_6,I_9,I_{12},I_{15},I_{18},I_{27},
J_9,J_{12},J_{15},J_{18},J_{21}^{a},J_{21}^{b}].
\tag{B61}
\] Moreover the fifty displayed secondary monomials are a free basis
over the polynomial ring in the seven parameters.

\textbf{Proof.} Work first over \(\mathbb Q\). Theorem 5.4 makes the
parameter ring \(P\) a polynomial ring. Lemma 6.2 and (B55)--(B60) make
the graded map \(\bigoplus_{W}P [-\deg W]\to R\), sending a basis
element to \(W\), injective. The table and Lemma 6.1 give exactly the
same Hilbert series for its domain and codomain. Each graded piece is
finite-dimensional; an injective map between spaces of equal finite
dimension is surjective. Therefore \(R=\bigoplus_W PW\). Every \(W\) is
a monomial in the six additional invariants, so (B61) follows.

For any characteristic-zero extension, the invariant spaces commute with
scalar extension degree by degree, by the finite rational derivation
kernels established in Theorem 5.4. The same is true of the displayed
free \(P\)-module. Extending the rational equality proves both
assertions over \(k\). Thus no external degree bound, unproved
regular-sequence assertion or specialization shortcut is needed.
\(\square\)

This proves the full thirteen-generator theorem and identifies a
concrete finite module basis. The certificate package includes the
all-degree character calculation as well as the finite independence
calculation, and its reproduction instructions keep the two logical
steps separate.

\subsection{Exercises}\label{exercises}

\protect\phantomsection\label{exercise-1}{}

\textbf{Exercise 1 --- The determinant character.} Show that a bracket
of \(n\) vector columns is special-linear invariant and changes by
\(\det g\) under \(v_i\mapsto gv_i\).

\hyperref[solution-1]{Solution}.

\protect\phantomsection\label{exercise-2}{}

\textbf{Exercise 2 --- Density from unit quaternions.} Prove the density
of \(\operatorname{SU}(2)\) directly from its quaternion
parametrization.

\hyperref[solution-2]{Solution}.

\protect\phantomsection\label{exercise-3}{}

\textbf{Exercise 3 --- Averaging a linear-form representation.} Compute
\(\rho\) on \(\mathbb C [x,y]_2\) for the standard special-linear
action, and determine it in every positive degree.

\hyperref[solution-3]{Solution}.

\protect\phantomsection\label{exercise-4}{}

\textbf{Exercise 4 --- Finite generation and normality.} Deduce that
\(B\) is a Noetherian normal domain. Explain why finite generation alone
does not prove normality, using \(\mathbb C [t^2,t^3]\).

\hyperref[solution-4]{Solution}.

\protect\phantomsection\label{exercise-5}{}

\textbf{Exercise 5 --- The cubic nullcone.} Determine the nullcone of
binary cubics and identify its defining invariant.

\hyperref[solution-5]{Solution}.

\subsection{Solutions}\label{solutions}

\subsubsection{Solution 1}\label{solution-1}

The new column matrix is \(g(v_{i_1}\ \cdots\ v_{i_n})\).
Multiplicativity of the determinant gives
\([gv_{i_1}\cdots gv_{i_n}]=\det(g)[v_{i_1}\cdots v_{i_n}]\). On
\(\operatorname{SL}_n\) the factor is one. Under the induced action on
polynomial functions, \(gF(v)=F(g^{-1}v)\), the character is its
inverse. Keeping these two assertions separate avoids a sign change in
determinant weights.

\hyperref[exercise-1]{Exercise 1}.

\subsubsection{Solution 2}\label{solution-2}

Every element is \[
\begin{pmatrix}
u+iv&s+it\\-s+it&u-iv
\end{pmatrix},
\qquad u,v,s,t\in\mathbb R,\quad u^2+v^2+s^2+t^2=1.
\] To check the parametrization, a unitary determinant-one matrix has
second row equal to the negative conjugate-swapped first row, and the
norm condition gives the sphere equation; the converse follows by
multiplication.

Near the identity use real coordinates \(v,s,t\), with
\(u=(1-v^2-s^2-t^2)^{1/2}\) and the positive root. The same formula is
holomorphic for complex \(v,s,t\) near zero, using the square root with
value one. Its derivative gives the three complex-independent matrices
\[
\begin{pmatrix}i&0\\0&-i\end{pmatrix},\quad
\begin{pmatrix}0&1\\-1&0\end{pmatrix},\quad
\begin{pmatrix}0&i\\i&0\end{pmatrix},
\] a basis of \(\mathfrak{sl}_2(\mathbb C)\). Thus it parametrizes a
complex open neighbourhood of the identity. A polynomial vanishing on
the real unit sphere vanishes on the real coordinate box here; the
coordinatewise identity theorem makes it vanish on that complex
neighbourhood. Connectedness, proved in
\href{NOE-INV-01.html\#lemma-8-2}{Lemma 8.2} by elementary matrices,
makes it vanish on all of \(\operatorname{SL}_2(\mathbb C)\).

\hyperref[exercise-2]{Exercise 2}.

\subsubsection{Solution 3}\label{solution-3}

Transitivity on nonzero vectors shows that an invariant polynomial on
\(\mathbb C^2\) is constant: it has the same value on the punctured
plane, a dense set. Hence the invariant subspace in every positive
homogeneous degree is zero. A Reynolds operator projects onto that
subspace, so \(\rho(ax^2+bxy+cy^2)=0\), and likewise in every positive
degree. On constants it is the identity. In particular \(\rho(x^2)\) is
not obtained by keeping only a torus-fixed monomial; being fixed by the
whole group is the condition.

\hyperref[exercise-3]{Exercise 3}.

\subsubsection{Solution 4}\label{solution-4}

\href{NOE-INV-01.html\#theorem-10-1}{Theorem 10.1} makes \(B\) a
quotient of a finite polynomial ring, hence Noetherian by the Hilbert
basis theorem. It is a domain because \(A\) is, and the monic-equation
argument of \href{NOE-INV-01.html\#corollary-10-2}{Corollary 10.2}
proves it is integrally closed in its fraction field. For
\(R=\mathbb C [t^2,t^3]\), the element \(t=t^3/t^2\) belongs to
\(\operatorname{Frac}R\), and satisfies the monic polynomial
\(Z^2-t^2\in R [Z]\). But \(t\notin R\): every nonconstant monomial in
its generators has degree at least two. Thus \(R\) is finite type and a
domain but is not normal.

\hyperref[exercise-4]{Exercise 4}.

\subsubsection{Solution 5}\label{solution-5}

\href{NOE-INV-01.html\#theorem-13-2}{Theorem 13.2} gives root
multiplicity \(>3/2\), hence at least two, for every nonzero cubic in
the nullcone. The zero cubic is also there. The first lesson proves that
the invariant ring is \(\mathbb C [\Delta]\), with
\(\Delta=(2/27)[12]^2 [13]^2 [23]^2\). Its zero set is exactly the
repeated-root locus, including zero, so \(\mathcal N=V(\Delta)\).
Alternatively, a double-root cubic can be moved to one divisible by
\(x^2\), and the diagonal one-parameter substitution in
\href{NOE-INV-01.html\#theorem-13-2}{Theorem 13.2} sends it to zero; a
three-distinct-root cubic has nonzero \(\Delta\).

\hyperref[exercise-5]{Exercise 5}.

\subsection{Further reading}\label{further-reading}

Noether's 1910 announcement and 1911 paper develop matrix products and
decomposition identities. Artebani and Dolgachev,
\href{https://arxiv.org/abs/math/0611590}{The Hesse pencil of plane
cubic curves}, treats the cubic geometry. The quartic constructions are
due to Dixmier and Ohno; Ohno's
\href{https://aeb.win.tue.nl/math/ohno-preprint.2007.05.15.pdf}{preprint},
Lercier, Liu, Lorenzo García and Ritzenthaler's
\href{https://arxiv.org/abs/1803.05816}{Reduction type of smooth plane
quartics}, and Elsenhans's
\href{https://math.uni-paderborn.de/fileadmin-eim/mathematik/AG-Computeralgebra/Publications-elsenhans/inv_pl_quart_5a.pdf}{Explicit
computations of invariants of plane quartic curves} describe the
parameters and generators.

\end{document}
